\capitulo{3}{Conceptos teóricos}

En este apartado considero importante definir que es un RAG y cuales son sus aportaciones más significativas a los LLM que ya existen, pero, también, considero importante contextualizar su origen, y para ello voy a empezar comentando que son los LLM y sus principales limitaciones. 

\section{LLM}
Los Large Language Model (LLM) son los grades modelos del lenguaje. Es decir, son modelos de aprendizaje automático que han sido entrenados con conjuntos de datos generados por humanos. Estos datos los usan para encontrar patrones estadísticos y replicar nuestras habilidades lingüísticas. Por lo cual podemos decir que en esencia son modelos de predicción del siguiente token o lo que es lo mismo de la siguiente palabra \citep{book:RAG_SimpleGuide}.

Algunos ejemplos de LLM son ChatGPT, Claude o Llama.

\subsection{Características}
- Naturaleza de predicción del próximo token: como he mencionado anteriormente los LLM seleccionan la palabra más probable de una distribución. La elección de esta palabra se basa en los patrones estadísticos que ha aprendido con los datos del entrenamiento.

- Memoria paramétrica: es el conocimiento interno del LLM. Consiste en la información con la que ha sido entrenado y se basa únicamente en sus parámetros. Por lo cual es una memoria limitada y estática.

\subsection{Limitaciones}
- Desactualizado: los eventos posteriores al entrenamiento del LLM, es decir, la fecha de corte del cocimiento, no estarán disponibles para el modelo lo que puede hacer que la respuesta este desactualizada. Además, el entrenamiento de un LLM es muy costoso y lleva mucho tiempo, por lo cual no se reentrenan con mucha frecuencia. 

- Alucinaciones: a veces los LLM dan respuestas incorrectas con gran confianza de manera que parece que la información que te está dando es verdadera. Esto se debe a su naturaleza de predicción del próximo token. 

- Limitación de conocimiento: los LLM se entrenan con datos públicos esto limita su conocimiento, ya que no tienen conocimientos de información que no sea publica, como, por ejemplo, documentos internos de una empresa, información de los clientes o de los productos.

\subsection{Solución}
Para solventar las limitaciones que hemos comentado la solución es darle al LLM el contexto que le falta. Para proporcionarle este contexto podemos expandir la memoria del LLM mediante generación aumentada, es decir, incorporando un RAG.

\section{RAG}
Retrieval Augmented Generation (RAG), o en español, generación aumentada por recuperación es la técnica que consiste en recuperar la información que sea relevante de una fuente externa, con dicha información aumenta la entrada al LLM, lo que permite que el LLM genere una respuesta mucho más precisa, contextual y confiable. Para realizar esto el RAG usa tres pasos, recuperar, aumentar y generar. 

EL RAG combina dos tipos de memoria:

- La memoria paramétrica: es la que épicamente tienen los LLM, que como hemos visto es limitada y estática. 

- La memoria no paramétrica: es información externa al LLM que se le puede proporcionar. Es un conocimiento externo y dinámico que se puede extender tanto como queramos y que puede contener documentos propietarios.

Por lo cual conociendo esto, podemos decir que la generación aumentada por recuperación (RAG) es una metodología que busca ampliar la memoria paramétrica de un LLM incorporando una memoria no paramétrica explícita. A través de un recuperador, se obtiene información relevante de esta memoria externa, la cual se añade al prompt y luego se envía al LLM. De esta forma, el modelo puede producir respuestas más contextuales, confiables y precisas \citep{book:RAG_SimpleGuide}.

\subsection{Objetivos del RAG}
- Hacer que los LLM respondan con información actualizada.

- Hacer que los LLM respondan información precisa, contextual y confiable.

- Hacer que los LLM conozcan información propietaria.

\subsection{Ventajas}
- Conocimiento del contexto profundo: esto se debe a que la información adicional ayuda al LLM a generar respuestas contextualmente apropiadas, ya que, las respuestas se basan en datos específicos. Esto aumenta la confianza del usuario.

- Citar fuentes: como la información se extrae de fuentes conocidas las puede citar en su respuesta. Esto dispara la fiabilidad y permite a los usuarios comprobar la información obtenida.

- Reduce las alucinaciones: el sistema esta anclado a los hechos, por lo que se reducen las respuestas inexactas o falsas. 

\section{Conexión LLM con RAG}
Cogemos los documentos externos y los convertimos en vectores numéricos (embeddings) y los almacenamos. Cuando llega una consulta, en lugar de buscar coincidencias por palabras buscamos en ese espacio vectorial los trozos que son semánticamente más parecidos, a esto se le conoce como búsqueda por significado. Esos fragmentos se incorporan como contexto al LLM, de modo que RAG combina la capacidad de razonamiento del modelo con una recuperación de información altamente precisa.

%Capitulo dos del libro
\section{Diseño del RAG}
Un sistema RAG se construye en dos pipelines complementarios, el de indexación, que es offline o asíncrono, y el de generación, que es online o en tiempo real. El primero crea y mantiene la memoria no paramétrica o base de conocimiento; el segundo gestiona la interacción con el usuario, realizando la recuperación, el aumento del prompt y la generación de la respuesta. Esta separación es estructural en el diseño de RAG y habilita tanto precisión como latencia controlada \citep{book:RAG_SimpleGuide}.

\subsection{Pipeline de Indexación (offline)}
Su objetivo es consolidar y preparar el conocimiento para que el pipeline de generación pueda consultarlo eficientemente. Consta de cinco pasos, el primero consiste en conectar a las fuentes externas, el segundo en extraer y parsear los documentos, el tercero en dividir textos largos en fragmentos más pequeños denominados chunks, el cuarto convierte esos chunks a un formato adecuado, típicamente los combatirte a embeddings vectoriales y el quinto almacena en una base de datos vectorial los embeddings. Además de crearlo, este pipeline mantiene y actualiza el conocimiento base, por lo que debe ejecutarse de forma periódica \citep{book:RAG_SimpleGuide}.

\subsubsection{Componentes clave}
\begin{itemize}
\tightlist
	\item Carga de datos (data loading): conectores a archivos o sistemas, por ejemplo, PDF, DOCX, JSON o HTML, extracción y parsing, preprocesado y gestión de metadatos.
	\item División de texto (Chunking): segmentación en unidades manejables, los chunks, para mejorar la recuperación y el alineamiento con las capacidades de contexto LLM.
	\item Conversión a vectores (Embeddings): representación numérica optimizada para búsqueda de similitud y operaciones de recuperación.
	\item Almacenamiento: bases de datos vectoriales para la búsqueda eficiente sobre embeddings y metadatos.
\end{itemize} 

Cuando la recuperación se externaliza, no siempre es necesario construir un pipeline de indexación propio \citep{book:RAG_SimpleGuide}. 

\subsection{Pipeline de Generación (online)}
Gestiona la consulta del usuario de extremo a extremo. Este pipeline recibe la pregunta del usuario, recupera el contexto relevante de la base vectorial, los chunks, aumenta el prompt con ese contexto y genera la respuesta con el LLM \citep{book:RAG_SimpleGuide}. 

\subsubsection{Componentes clave}
\begin{itemize}
\tightlist
	\item Retriever: es el motor de búsqueda sobre la base vectorial. Es crítico para la precisión del sistema y contribuye a la latencia, su calidad condiciona el rendimiento del RAG.
	\item Gestión del Prompt: combina la consulta y el contexto recuperado. Además, aplica estrategias de ingeniaría de prompt para maximizar la calidad de la respuesta.
	\item Configuración LLM: es la selección del modelo que va a generar la respuesta final. Pueden ser modelos fundacionales, fine-tuned, abiertos o de servicio gestionado.
\end{itemize} 

\subsection{Capas de soporte}
Además de los dos pipelines, un sistema en producción requiere orquestación, evaluación y monitorización continua, cache semántico para reducir coste y latencia, controles de seguridad para cumplimiento normativo y seguridad frente a inyecciones en el prompt, envenenamiento de dato o fugas de información. Todo ello se sostiene sobre una infraestructura de servicio y un RAGOps stack con capas de datos, modelos, prompting, evaluación, orquestación, despliegue, hosting y monitorización \citep{book:RAG_SimpleGuide}.

\section{Evaluación de calidad del RAG}
Se hará mediante la triada de evaluación propuesta por TruEra. Estructura la calidad en tres comprobaciones entre consulta (prompt), contexto recuperado y respuesta. Complementariamente, se emplean métricas como relevancia, precisión y recall, y conjuntos ground truth para validación objetiva y monitorización continua en producción \citep{book:RAG_SimpleGuide}.

\subsection{Las tres dimensiones de la triada:}
\begin{itemize}
\tightlist
	\item Relevancia del contexto: la información que hemos recuperado tiene que ver con la pregunta.
	\item Groundedness o fidelidad: la respuesta que da el LLM se basa de verdad en el contexto que le hemos dado, sin alucinaciones ni omisiones clave.
	\item Relevancia de la respuesta: la respuesta es útil y adecuada respecto a la intención y el contexto de la consulta original.
\end{itemize}